\section{Discussion}\label{sec:discussion}
Purification preserves the reduced density matrix but can change the
spectral measure defining complexity. With the source's normalized seeds,
each side of the mixed/purification hierarchy fails in an exact full-rank
qutrit example. The qubit theorem fixes the minimum physical dimension
for both failures. The upper-side perturbation further shows that its
failure does not require an exactly decoupled spectator.

The ratio results delimit a different issue. Uniform temporal
concentration of $r$ and of its reciprocal fail by different four-level
families, and fixed purity cannot supply finite universal upper control
of $r$ or its mean. These statements concern explicit universal
formalizations of the qualitative source discussion. They neither replace
that discussion by an invented quantitative conjecture nor invalidate its
sample-specific numerical observations. Typical-state statements,
restricted spectral families, and assumptions controlling the smallest
populations would require their own quantified formulations.

The surviving factor-two inequality and qubit bounds demonstrate that
useful comparisons remain possible when their spectral structure is
specified. The pure-density inequality belongs to prior tensor-product
superadditivity and has been included to separate that protected
comparison from the failed mixed-state ordering. The state-dependent
return bounds identify information that a replacement comparison can
retain. No universal classification of such replacements is asserted.

All decisive constructions are finite dimensional. The temporal
minimum-dimension question between three and four remains open within
this analysis, and no claim is made about unrestricted infinite-dimensional
or nonunitary dynamics. The accompanying exact certificates make the
stated finite results directly checkable without treating numerical plots
as proofs.
